%****************************************************
%	CHAPTER 9 - CONCLUSION
%****************************************************
\chapter{Conclusion and Recommendations}
\label{ch:conclusion}
%====================================================
\section{Conclusion}
\label{sec:ch9.section1}
%----------------------------------------------------

This investigation set out to explore the use of an off-the-shelf \ac{LiDAR} sensor, the Livox Avia, for remote sensing of the Antarctic \ac{MIZ}, with a focus on assessing the sensor's capability to produce usable point cloud data for the estimation of various sea ice field parameters. The primary objectives were to first characterise the sensor's performance in controlled settings to obtain a further understanding of its capabilities and limitations, before moving on to the processing and analysis of point cloud data collected from the Aalto University Wave and Ice Tank and finally in-situ data collected in the \ac{MIZ} during the \ac{SCALE} Winter 2022 cruise.

Through engaging with the relevant literature surrounding the \ac{MIZ}, sea ice found in the region (particularly pancake ice), the use of \ac{LiDAR} as a remote sensing tool and common methods to process and analyse point cloud data, parameters of interest were identified for estimation from the \ac{LiDAR} data, as well as potential methods to extract these parameters. The parameters selected were surface wave characteristics (wave amplitude and period), and pancake ice floe identification. To extract these parameters from the \ac{LiDAR} point cloud data, a processing pipeline was proposed, consisting a pre-processing stage to compensate for motion of the observation platform through the use of a Kalman Filter-based approach, followed by a de-noising filter to remove invalid points in the point cloud data. The main processing stage then consisted of two separate branches, each consisting of specific feature extraction methods for the parameters of interest. A \ac{RANSAC} based plane fitting method was proposed for surface wave parameter estimation, while a density-based clustering stage was proposed for pancake ice floe identification.

During the sensor characterisation phase of this investigation, it was found that the Livox Avia struggled to accurately reconstruct simple geometric targets under conditions simulating those expected during Antarctic deployments. A combination of observing targets at acute angles of incidence, the sensor's relatively large beam divergence and low point densities led to distortions in the reconstructed targets and their estimated dimensions. Additionally, it was found that attempts to estimate the point normals in the targets did not yield any further insights over the raw point cloud data, due to the low point densities and distortions present in the scenes. Despite these limitations, the sensor was still able to produce point cloud data from which approximate target dimensions could be estimated, and rough geometries could be inferred, indicating that the sensor shows promise for medium-scale environmental observations.

The pre-processing stage of the proposed processing pipeline was found to be effective in reducing motion artefacts present in the \ac{LiDAR} point cloud data collected during the \ac{SCALE} cruise. While the \ac{ESKF} implementation was limited to compensating for pitch and roll motions only, significant improvements in point cloud quality were observed following motion compensation. The de-noising filter was also found to be effective in removing invalid points from the point cloud data, however, it was noted that this stage is not necessarily required if the data being processed is already of high quality.

The surface wave parameter estimation branch of the processing pipeline produced positive results, with estimated wave periods matching the ground truth data with a high degree of accuracy. Wave amplitude estimations were found to be slightly less accurate, with a tendency to overestimate wave amplitudes. Overall the wave estimation method proved highly effective in controlled wave tank settings, demonstrating the potential for \ac{LiDAR}-based wave observations in ice-covered waters provided an appropriate wave model is used.

The pancake ice floe identification branch of the processing pipeline produced mixed results. All three clustering algorithms tested struggled to consistently identify entire pancake floes in the point cloud data across the various scenes analysed. Generally, the algorithms were successful in identifying the edges of pancake floes, where point densities were higher due to the ridged geometries present at pancake boundaries, but struggled to identify the interior regions of the pancakes, where point densities were lower. \ac{DBScan} and \ac{OPTICS} were found to perform acceptably with careful parameter tuning, while the implementation of \ac{VDBScan} used in this investigation had a notable failure and thus is not recommended for further use in its current form. Overall, while the \ac{DBScan} and \ac{OPTICS} algorithms showed some promise for pancake ice floe identification, further work is required to improve their performance and reliability.

An aspect of the \ac{LiDAR} data that was underutilised during this investigation was the point return intensities. It was observed that the point intensities varied across individual floes and ice types, with higher intensities often observed at pancake edges and ridges. However, these intensity variations were not incorporated into any stages of the processing pipeline. It is likely that further insights are to be gained from a deeper analysis of the point intensities present in the data.

In conclusion, the Livox Avia sensor demonstrated potential as a remote sensing tool for data collection in the Antarctic \ac{MIZ}. Despite the identified limitations of the sensor hardware, the point cloud data produced was of sufficient quality to enable the extraction of useful environmental parameters. Further work in refining the processing pipeline is recommended but is a promising step towards utilising \ac{LiDAR} technology for Antarctic sea ice observations. Overall, this investigation has highlighted the potential of \ac{LiDAR} viable tool for Antarctic \ac{MIZ} research.

%====================================================
\section{Recommendations for Future Work}
\label{sec:ch9.section2}
%----------------------------------------------------

Both Antarctic research, particularly observation of the \ac{MIZ} and the use of \ac{LiDAR} as a remote sensing tool are emerging fields with lots of scope for future work. Some recommendations for future work building on this project include:

A detailed analysis of the Livox Avia's beam divergence and its scene reconstruction accuracy under both ideal targets and simulated environmental conditions. Conducting a thorough characterisation of the sensor's performance in controlled settings, replicating expected deployment temperatures and using simple targets at a 0° angle of incidence, could provide further insight into the sensor's capabilities and limitations for Antarctic deployments.

The addition of additional motion sensors synced to the Avia, such as a magnetometer or \ac{GPS} unit, to assist in more thorough motion compensation of the \ac{LiDAR} point cloud data. This would assist in further reducing motion artefacts in the data induced by yaw and translation of the observation platform.

An investigation into more advanced wave models that may better capture the behaviour of surface waves propagating through ice-covered waters. This could allow for the estimation of wave parameters from more complex wave fields found in the \ac{MIZ}, such as those observed during the \ac{SCALE} cruise.

An investigation into the performance of custom or modified clustering algorithms for improved pancake ice floe identification. This could include the development of clustering algorithms that incorporate additional features beyond spatial coordinates, such as intensity values. An analysis of the density and distribution of floes present in the point cloud data could also be conducted to better inform algorithm parameter tuning and design.

A deeper analysis into the \ac{LiDAR} point return intensities and their relationship to different ice types or surface parameters. The variation in point return intensities present in the data collected during this investigation suggests that there may be useful information not yet extracted from the data. A study into how these intensities vary across different ice types (frazil, nilas, pack, etc.) or surface parameters (roughness, geometry, angle of incidence, etc.) are likely to be of value. Insights from such a study could inform future clustering algorithm designs, ice type classification efforts or surface parameter estimation.

The use of multiple optical sensors to complement the \ac{LiDAR} data collection. The addition of sensors such as thermal cameras or video/stereo cameras to form a multi-sensor suite would provide additional data streams to enhance observations of the \ac{MIZ}. Specifically in the identification and estimation of ice floe geometry and dimensions, thermal sensors could assist in distinguishing between individual floes or ice types based on temperature differences, while video or stereo cameras could provide visual context to the \ac{LiDAR} data, aiding in ice type classification and surface feature identification.

%****************************************************
% END
%****************************************************